% ---- score2 standard preamble (from: python3 score2_tools.py preamble) ----
\documentclass[10pt]{extarticle}
\usepackage[a4paper,margin=1.5cm]{geometry}
\usepackage[T1]{fontenc}
\usepackage{amsmath,amssymb}
\usepackage{graphicx}
\usepackage[export]{adjustbox}
\usepackage{multicol}
\usepackage{enumitem}
\usepackage{booktabs,array,tabularx}
\setlength{\columnsep}{0.7cm}
\setlength{\parindent}{0pt}
\setlength{\parskip}{3pt}
\pagestyle{empty}
\makeatletter
% figure registry (figdefs.tex is generated by the extract step)
\newcommand{\qfigdef}[4]{\@namedef{qf@file@#1}{#2}\@namedef{qf@w@#1}{#3}\@namedef{qf@wide@#1}{#4}}
\newcommand{\qf@check}[1]{\@ifundefined{qf@file@#1}{\PackageError{score2}{Unknown figure id #1}{}}{}}
\newcommand{\qf@img}[4][]{\includegraphics[width=#2\dimexpr\@nameuse{qf@w@#4}pt\relax,max width=#3,max height=0.45\textheight,#1]{\@nameuse{qf@file@#4}}}
% Full-width questions: finalize wraps every question that contains a wide figure in \qfullbegin...\qfullend,
% so the WHOLE question (stem, figures, options) is set across the page in reading order, then the two columns
% resume. The question's head - its number, stem and first figure (\qfullsplit marks the end of the head) -
% is kept on one page, so a number is never left at a page foot with its figure overleaf; the rest may break
% across pages normally, so no big gaps are left. Without a \qfullsplit, the whole question is kept together
% if it is under 45% of a page. Never write these yourself.
\newif\ifqf@full
\newif\ifqf@boxed
\newsavebox{\qf@box}
\newcommand{\qf@open}{\global\qf@fulltrue\par\global\qf@boxedtrue
  \setbox\qf@box\vbox\bgroup\hsize\textwidth\linewidth\textwidth}
\newcommand{\qf@emit}[1]{\par\egroup\global\qf@boxedfalse
  \ifdim\dimexpr\ht\qf@box+\dp\qf@box\relax<#1\textheight\noindent\usebox{\qf@box}\par\else\unvbox\qf@box\fi}
\newcommand{\qfullsplit}{\ifqf@boxed\qf@emit{0.85}\fi}
\newcommand{\qf@close}{\ifqf@boxed\qf@emit{0.45}\fi\global\qf@fullfalse}
\newcommand{\qfullbegin}{\end{multicols}\qf@open}
\newcommand{\qfullend}{\qf@close\begin{multicols}{2}}
\newcommand{\qfullstart}{\qf@open}
\newcommand{\qfulldone}{\qf@close}
% \qfig[scale]{id}: block figure, centred, directly where it is written (size set by the tool).
% Fallback if a wide figure is not inside a full-width question: break out of the columns in place - no floats.
\newcommand{\qfig}[2][1]{\qf@check{#2}%
  \ifqf@full\qf@blk{#1}{#2}\else
  \ifnum\@nameuse{qf@wide@#2}=1 \expandafter\@firstoftwo\else\expandafter\@secondoftwo\fi
  {\end{multicols}\begin{center}\qf@img{#1}{\textwidth}{#2}\end{center}\begin{multicols}{2}}%
  {\qf@blk{#1}{#2}}\fi}
\newcommand{\qf@blk}[2]{\par\smallskip{\centering\qf@img{#1}{\linewidth}{#2}\par}\smallskip}
% \qfigi[scale]{id}: inline figure (inside option lists / tables / minipages); never breaks columns.
\newcommand{\qfigi}[2][1]{\qf@check{#2}\qf@img[valign=c]{#1}{\linewidth}{#2}}
\makeatother
% structure macros - lint reads these to check completeness and figure placement
\newcommand{\papersec}[1]{\par\bigskip{\centering\Large\bfseries PAPER-#1\par}\medskip}
\newcommand{\partsec}[2]{\par\medskip\noindent{\large\bfseries PART-#1: #2}\par\smallskip}
\newcommand{\sect}[1]{\par\smallskip\noindent\textbf{#1}\par}
\newcommand{\qn}[1]{\par\smallskip\noindent\textbf{#1.}~}
\newcommand{\paraq}[2]{\par\smallskip\noindent\textbf{Paragraph for Questions #1 and #2}\par}
\newenvironment{qoptions}{\begin{itemize}[leftmargin=1.8em,itemsep=1pt,topsep=2pt,label={}]}{\end{itemize}}
\newcommand{\opt}[2]{\item[(#1)] #2}
% figure registry (generated by score2_tools.py finalize)
\qfigdef{p1_phy_s1_q01_1}{Score-II - Test 26 (20 Mar 2024)-img/p1_phy_s1_q01_1.png}{155.0}{0}
\qfigdef{p1_phy_s1_q04_1}{Score-II - Test 26 (20 Mar 2024)-img/p1_phy_s1_q04_1.png}{54.8}{0}
\qfigdef{p1_phy_s2_q06_1}{Score-II - Test 26 (20 Mar 2024)-img/p1_phy_s2_q06_1.png}{108.3}{0}
\qfigdef{p1_phy_s2_q07_1}{Score-II - Test 26 (20 Mar 2024)-img/p1_phy_s2_q07_1.png}{108.0}{0}
\qfigdef{p1_phy_s2_q09_1}{Score-II - Test 26 (20 Mar 2024)-img/p1_phy_s2_q09_1.png}{102.4}{0}
\qfigdef{p1_phy_s3_para3_4_1}{Score-II - Test 26 (20 Mar 2024)-img/p1_phy_s3_para3_4_1.png}{79.5}{0}
\qfigdef{p1_phy_s3_para5_6_1}{Score-II - Test 26 (20 Mar 2024)-img/p1_phy_s3_para5_6_1.png}{117.4}{0}
\qfigdef{p1_chem_s1_q04_A}{Score-II - Test 26 (20 Mar 2024)-img/p1_chem_s1_q04_A.png}{105.3}{0}
\qfigdef{p1_chem_s1_q04_B}{Score-II - Test 26 (20 Mar 2024)-img/p1_chem_s1_q04_B.png}{103.8}{0}
\qfigdef{p1_chem_s1_q04_C}{Score-II - Test 26 (20 Mar 2024)-img/p1_chem_s1_q04_C.png}{103.6}{0}
\qfigdef{p1_chem_s1_q04_4}{Score-II - Test 26 (20 Mar 2024)-img/p1_chem_s1_q04_4.png}{104.1}{0}
\qfigdef{p1_chem_s2_q05_A}{Score-II - Test 26 (20 Mar 2024)-img/p1_chem_s2_q05_A.png}{127.6}{0}
\qfigdef{p1_chem_s2_q05_B}{Score-II - Test 26 (20 Mar 2024)-img/p1_chem_s2_q05_B.png}{120.8}{0}
\qfigdef{p1_chem_s2_q05_C}{Score-II - Test 26 (20 Mar 2024)-img/p1_chem_s2_q05_C.png}{132.5}{0}
\qfigdef{p1_chem_s2_q05_4}{Score-II - Test 26 (20 Mar 2024)-img/p1_chem_s2_q05_4.png}{85.0}{0}
\qfigdef{p1_chem_s2_q06_A}{Score-II - Test 26 (20 Mar 2024)-img/p1_chem_s2_q06_A.png}{109.4}{0}
\qfigdef{p1_chem_s2_q06_B}{Score-II - Test 26 (20 Mar 2024)-img/p1_chem_s2_q06_B.png}{128.8}{0}
\qfigdef{p1_chem_s2_q06_C}{Score-II - Test 26 (20 Mar 2024)-img/p1_chem_s2_q06_C.png}{160.9}{0}
\qfigdef{p1_chem_s2_q06_4}{Score-II - Test 26 (20 Mar 2024)-img/p1_chem_s2_q06_4.png}{209.1}{0}
\qfigdef{p1_chem_s2_q08_1}{Score-II - Test 26 (20 Mar 2024)-img/p1_chem_s2_q08_1.png}{130.3}{0}
\qfigdef{p1_chem_s2_q09_1}{Score-II - Test 26 (20 Mar 2024)-img/p1_chem_s2_q09_1.png}{203.5}{0}
\qfigdef{p1_chem_s3_para1_2_1}{Score-II - Test 26 (20 Mar 2024)-img/p1_chem_s3_para1_2_1.png}{43.9}{0}
\qfigdef{p1_chem_s3_q01_1}{Score-II - Test 26 (20 Mar 2024)-img/p1_chem_s3_q01_1.png}{194.9}{0}
\qfigdef{p1_chem_s4_q07_1}{Score-II - Test 26 (20 Mar 2024)-img/p1_chem_s4_q07_1.png}{271.3}{1}
\qfigdef{p1_chem_s4_q07_2}{Score-II - Test 26 (20 Mar 2024)-img/p1_chem_s4_q07_2.png}{219.4}{0}
\qfigdef{p1_chem_s4_q07_3}{Score-II - Test 26 (20 Mar 2024)-img/p1_chem_s4_q07_3.png}{52.7}{0}
\qfigdef{p1_chem_s4_q07_4}{Score-II - Test 26 (20 Mar 2024)-img/p1_chem_s4_q07_4.png}{74.9}{0}
\qfigdef{p1_chem_s4_q07_5}{Score-II - Test 26 (20 Mar 2024)-img/p1_chem_s4_q07_5.png}{211.1}{0}
\qfigdef{p1_chem_s4_q07_6}{Score-II - Test 26 (20 Mar 2024)-img/p1_chem_s4_q07_6.png}{128.6}{0}
\qfigdef{p1_chem_s4_q08_1}{Score-II - Test 26 (20 Mar 2024)-img/p1_chem_s4_q08_1.png}{47.0}{0}
% ---- body skeleton ----
% \begin{document}
% \begin{center}{\Large\bfseries JEE (Main+Advanced) Score-II --- Test NN (DD Mon YYYY)}\end{center}
% \begin{multicols}{2}
%   \papersec{1}   (only for two-paper tests)
%   \partsec{1}{PHYSICS} \sect{Section-I (i): Single Correct} \qn{1} ... \qfig{p1_phy_q01_1} ...
% \end{multicols}
% \end{document}

\begin{document}
\begin{center}{\Large\bfseries JEE (Main+Advanced) Score-II --- Test 26 (20 Mar 2024)}\end{center}
\begin{multicols}{2}

\partsec{1}{PHYSICS}
\sect{Section-I(i)}

\qn{1} If the figure shown, a box (1), a wedge (2) and block (3) are arranged as shown on a frictionless surface. The angle of inclination of the surface of wedge is $\theta=\sin^{-1}\left(\dfrac35\right)$. The masses of 1, 2 \& 3 are same. The accelerations of 1 \& 2 are respectively a \& c and that of 3 relative to 2 is b. The relation between a, b \& c is
\qfig{p1_phy_s1_q01_1}
\begin{qoptions}
\opt{A}{5a + 4b = 10c}
\opt{B}{4a + 5b = 10c}
\opt{C}{3a + 5b = 4c}
\opt{D}{3a + 4b = 5c}
\end{qoptions}

\qn{2} The workdone on a particle of mass m by a force $\vec F=k\left(\dfrac{x}{(x^2+y^2)^{3/2}}\hat\imath+\dfrac{y}{(x^2+y^2)^{3/2}}\hat\jmath\right)$ N (k is a constant of appropriate dimensions), when the particle is taken from a point (a, 0) to $\left(\sqrt3a,a\right)$ along a straight line path is
\begin{qoptions}
\opt{A}{$\dfrac ka\ln2$ J}
\opt{B}{$\dfrac{k}{2a}\ln2$ J}
\opt{C}{$\dfrac{k}{2a}$ J}
\opt{D}{$\dfrac ka$ J}
\end{qoptions}

\qn{3} Two moles of an ideal monoatomic gas undergo a thermodynamic process in which the molar heat capacity `C' of the gas depends on absolute temperature as $C=\dfrac{RT}{T_0}$ where R is a gas constant and T$_0$ is the initial temperature of the gas. (V$_0$ is the initial volume of the gas). Mark the correct statement about the gas process.
\begin{qoptions}
\opt{A}{Volume is minimum at T = T$_0$.}
\opt{B}{Volume is minimum at T = $\dfrac{T_0}{2}$.}
\opt{C}{Volume is minimum at T = $\dfrac32T_0$.}
\opt{D}{Volume is minimum at T = 2T$_0$.}
\end{qoptions}

\qn{4} Two uniform square plates each of mass m and side $\sqrt2\ell$ are connected by a uniform rod of length $\ell$ and mass m as shown. Time period of SHM in the plane of figure about point P is given by
\qfig{p1_phy_s1_q04_1}
\begin{qoptions}
\opt{A}{$2\pi\sqrt{\dfrac{8\ell}{g}}$}
\opt{B}{$2\pi\sqrt{\dfrac{12\ell}{g}}$}
\opt{C}{$2\pi\sqrt{\dfrac{7\ell}{3g}}$}
\opt{D}{$2\pi\sqrt{\dfrac{8\ell}{3g}}$}
\end{qoptions}

\sect{Section-I(ii)}

\qn{5} A particle moves in space such that its position $\vec r$ varies with time as\\
$\vec r=(ut+A\sin\omega t)\hat\imath+B\cos\omega t\,\hat\jmath+C\hat k$\\
The instant(s) at which the average velocity of the body is equal to u is/are
\begin{qoptions}
\opt{A}{$\dfrac\pi\omega$}
\opt{B}{$\dfrac{2\pi}\omega$}
\opt{C}{$\dfrac{3\pi}\omega$}
\opt{D}{$\dfrac{4\pi}\omega$}
\end{qoptions}

\qn{6} The given figure shows a circuit having an inductor, a capacitor and a resistor. Inductance L = 100 mH, capacitance C = 10$\mu$F and resistance $100\Omega$. The rms value of a.c. source is 200 V and angular frequency of source is 2000 rad/s.
\qfig{p1_phy_s2_q06_1}
\begin{qoptions}
\opt{A}{rms voltage across the inductor is $\dfrac{200}{\sqrt{13}}$ V.}
\opt{B}{rms voltage across the capacitor is $\dfrac{400}{\sqrt{13}}$ V.}
\opt{C}{rms voltage across the resistor is $\dfrac{600}{\sqrt{13}}$ V.}
\opt{D}{rms current through capacitor branch is $\dfrac{6}{\sqrt{13}}$ A.}
\end{qoptions}

\qn{7} Two discs A and B of same mass and radius R, rotating in opposite sense are joined with light rod having frictionless axels is placed on rough horizontal surface as shown in figure. The coefficient of kinetic friction between disc A and surface as well as B and surface is $\mu$. Initial velocity of each disc is zero. Then :-
\qfig{p1_phy_s2_q07_1}
\begin{qoptions}
\opt{A}{Final speed of B is $\dfrac{\omega_0R}{6}$.}
\opt{B}{Time after which rolling starts in A is $\dfrac{\omega_0R}{2\mu g}$.}
\opt{C}{Time after which rolling starts in B is $\dfrac{11\omega_0R}{12\mu g}$.}
\opt{D}{Maximum velocity of B is $\dfrac{\omega_0R}{4}$.}
\end{qoptions}

\qn{8} A long straight stationary wire is parallel to y-axis and passes through the point z = h on the z-axis. A current $I(t)=I_0\dfrac tT$ flows in the wire.\\
Lying in the x--y plane is a thin rectangular loop with two of its sides of length l, parallel to the long wire. The length b of other two sides is very small. The loop is located at x = x$_0$ at a given instant and its resistance is R, then :-
\begin{qoptions}
\opt{A}{At t = T and for $x_0=h\sqrt3$, the emf induced in the given loop is $e=\ell b\dfrac{\mu_0I_0\sqrt3}{8\pi hT}$}
\opt{B}{At t = T and for $x_0=h\sqrt3$, the torque acting on the loop is $\left(\dfrac{\mu_0I_0\ell b}{8\pi h}\right)^2\dfrac{\sqrt3}{TR}$.}
\opt{C}{At t = T and for x$_0$ = 0, the emf induced in the given loop is zero}
\opt{D}{At t = 0 and for x$_0$ = h, the torque acting on the loop is zero}
\end{qoptions}

\qn{9} A potentiometer wire PQ is made by using two wires of different cross-sectional area A and 2A as shown. The material of wire PS and SQ is same and their lengths are 5 m each. An ideal source battery of 10 V is connected in the circuit choose the correct option(s) regarding balance point.
\qfig{p1_phy_s2_q09_1}
\begin{qoptions}
\opt{A}{If the balance point is at a distance of 3 m from end P then emf e$_0$ = 4V}
\opt{B}{If the balance point is at a distance of 7m from end P then emf e$_0$ = 8V}
\opt{C}{If emf e$_0$ = 9V then balance point from end P is at 6.75 m.}
\opt{D}{Balance point is independent of internal resistance of the cell e$_0$.}
\end{qoptions}

\qn{10} A 100 eV electron collides with a stationary helium ion (He$^+$) in its ground state and excites it to a higher level. After the collision He$^+$ ions emit two photons in succession with wavelength 1084 \AA{} and 304\AA{} (Take hc = 12400 eV.\,\AA)
\begin{qoptions}
\opt{A}{The principal quantum number of excited state of He$^+$ due to collision is n = 3.}
\opt{B}{The shorter wavelength corresponds to transition from first excited state to ground state.}
\opt{C}{The recoil energy of the electron after collision is close to 48eV.}
\opt{D}{If the electron had a kinetic energy of 40 eV the recoil energy of electron would have been 40 eV.}
\end{qoptions}

\sect{Section-II(i)}

\paraq{1}{2}
Two soap bubbles A and B are kept in a closed chamber where the air is maintained at pressure 12 N/m$^2$. The radii of bubbles A and B is 2 cm and 4 cm respectively. Surface tension of the soap water used to make bubbles is 0.04 Nm$^{-1}$.

\qn{1} Find the ratio $\dfrac{n_B}{n_A}$, where n$_A$ and n$_B$ are the number of moles of air in bubbles A and B respectively.

\qn{2} The pressure in the chamber is now reduced to zero. The ratio $\left(\dfrac{r_B}{r_A}\right)^2$, where r$_B$ and r$_A$ are the new radii of the bubbles A and B, is

\paraq{3}{4}
The given circuit shown involves two ideal cells connected to a 1$\mu$F capacitor via a key K. Initially the key K is in position 1 and the capacitor is charged fully by 2V cell. The key is pushed to position 2.
\qfig{p1_phy_s3_para3_4_1}

\qn{3} After switching the key to position 2, the heat produced in $2\Omega$ resistor in a long time is \underline{\hspace{1cm}}$\mu$J.

\qn{4} After switching the key to position 2, the heat produced in $2\Omega$ resistor in $1\mu$s is \underline{\hspace{1cm}}

\paraq{5}{6}
Young's double slit experiment is conducted in a liquid of refractive index $\mu_1=\dfrac43$ as shown in figure. The separation S$_1$S$_2$ = d = 1.0 mm and D = 2 m. The wavelength of light incident is $\lambda$ = 5000\AA{} in air.
\qfig{p1_phy_s3_para5_6_1}

\qn{5} A thin transparent film of refractive index $\mu_2=\dfrac32$ is now pasted over slit S$_2$. The minimum thickness of the film so that the intensity at `O' becomes zero is \underline{\hspace{1cm}}$\mu$m.

\qn{6} Now further with the same arrangement as above, the minimum distance from the centre `O' of the screen where intensity is $\left(\dfrac14\right)^{\text{th}}$ of the maximum intensity is \underline{\hspace{1cm}}$\times10^{-4}$ m.

\sect{Section-II(ii)}

\qn{7} Electric field in a region of space is given as E = E$_0$\,x(x$-$1)(x$-$2)(x$-$3) N/C along x-axis, where E$_0$ is a positive constant. A positive charge `q' of mass m is in stable equilibrium at some specific position x, where x $>$ 0. If this charge is displaced slightly from equilibrium position to perform periodic motion, the angular frequency in rad/s of periodic motion will be [Given qE$_0$ = 2m]

\qn{8} A metal is heated in a furnace where a sensor is kept over the metal surface to measure the power radiated (P) by the metal. The sensor has a scale that displays $\log_2\left(\dfrac{P}{P_0}\right)$. Where P$_0$ is a constant. When the metal surface is at a temperature of 107$^\circ$C, the sensor shows zero. Assuming the emissivity of metallic surface remains constant, if the value displayed by the sensor is (10+n) when temperature is raised to 2767$^\circ$C, find n :

\qn{9} A 150W light bulb emits 5\% of its energy as electro magnetic radiation. The magnitude of average poynting S vector 10 m from the bulb is $\alpha\times10^{-3}$ W/m$^2$. Find $\alpha$. (Round off to nearest integer.)

\partsec{2}{CHEMISTRY}
\sect{Section-I(i)}

\qn{1} Total number of ether structure possible for molecular formula C$_5$H$_{12}$O : \fbox{$N_1$}\\
Total number of structure possible for molecular formula C$_4$H$_9$Cl : \fbox{$N_2$}\\
value of (N$_1$ + N$_2$) =
\begin{qoptions}
\opt{A}{6}
\opt{B}{8}
\opt{C}{10}
\opt{D}{9}
\end{qoptions}

\qn{2} During hydrolysis of compound CH$_3$--S--CH$_2$--CD$_2$--Cl two products are obtained. One product is CH$_3$--S--CH$_2$--CD$_2$--OH \& other product is :
\begin{qoptions}
\opt{A}{HO--CH$_2$--CD$_2$--OH}
\opt{B}{CH$_3$--S--CD$_2$--CH$_2$--OH}
\opt{C}{HO--S--CH$_2$--CD$_2$--OH}
\opt{D}{CD$_3$--S--CH$_2$--CD$_2$--OH}
\end{qoptions}

\qn{3} Choose the correct statements
\begin{qoptions}
\opt{A}{Cl$_2$ is strongest oxidising halogen since Cl$_2$ has highest EA}
\opt{B}{BDE order : F$_2$ $>$ Cl$_2$ $>$ Br$_2$ $>$ I$_2$}
\opt{C}{Maximum number of atoms in a plane for I$_2$Cl$_6$ and B$_2$H$_6$ are same.}
\opt{D}{BF$_3$ and BrF$_3$ both are planar.}
\end{qoptions}

\qn{4} A buffer solution is prepared by mixing `X' moles of CH$_3$COONa and `y' moles of CH$_3$COOH such that (x + y) = 1, in water to make 1 L buffer solution. If the buffer capacity of this buffer solution is plotted against moles of salt CH$_3$COONa (with x) then the plot obtained will be (to scale approximately).
\begin{qoptions}
\opt{A}{\qfigi{p1_chem_s1_q04_A}}
\opt{B}{\qfigi{p1_chem_s1_q04_B}}
\opt{C}{\qfigi{p1_chem_s1_q04_C}}
\opt{D}{\qfigi{p1_chem_s1_q04_4}}
\end{qoptions}

\sect{Section-I(ii)}

\qn{5} Which of following is/are configurational diastreomers :
\begin{qoptions}
\opt{A}{\qfigi{p1_chem_s2_q05_A}}
\opt{B}{\qfigi{p1_chem_s2_q05_B}}
\opt{C}{\qfigi{p1_chem_s2_q05_C}}
\opt{D}{\qfigi{p1_chem_s2_q05_4}}
\end{qoptions}

\qn{6} Correct option amongst the following is/are :
\begin{qoptions}
\opt{A}{\qfigi{p1_chem_s2_q06_A}}
\opt{B}{\qfigi{p1_chem_s2_q06_B}}
\opt{C}{\qfigi{p1_chem_s2_q06_C}}
\opt{D}{\qfigi{p1_chem_s2_q06_4}}
\end{qoptions}

\qn{7} Which of the following Statement is/are correct?
\begin{qoptions}
\opt{A}{Bis(benzoylacetonato)beryllium(II) can show optical isomerism}
\opt{B}{Square planar complexes are more favored sterically than tetrahedral complexes}
\opt{C}{C-O bond strength order : [Ni(CO)$_4$] $>$ [Co(CO)$_4$]$^-$ $>$ [Fe(CO)$_4$]$^{2-}$}
\opt{D}{[Zn(Cl)$_4$]$^{2-}$ is tetrahedral while [Zn(CN)$_4$]$^{2-}$ is square planar}
\end{qoptions}

\qn{8}
\qfig{p1_chem_s2_q08_1}
`B' + Ca(OH)$_2$(aq) $\longrightarrow$ solution turns milky\\
`B' gives green solution with K$_2$Cr$_2$O$_7$(aq) in acidic medium.\\
Find correct statement(s)
\begin{qoptions}
\opt{A}{`C' has structure similar to H$_2$O$_2$}
\opt{B}{`B' is SO$_3$}
\opt{C}{`E' forms a complex with CuSO$_4$(aq) solution}
\opt{D}{`D' is dibasic acid}
\end{qoptions}

\qn{9}
\qfig{p1_chem_s2_q09_1}
In the above following cell :\\
Pt $|$ H$_2$(g) (1 bar) $|$ NaA(aq) (0.1 M), HA(aq.) (0.8 M) $\|$ HCl (aq.) (0.1M) $|$ H$_2$(g) (1 bar) $|$ Pt\\
If 20 ml of 0.5 M NaOH(aq.) is added to anodic compartment then correct option(s) is/are :\\
(Given : log 5 = 0.7, log 2 = 0.3, $\dfrac{2.303\,RT}{F}=0.06$)
\begin{qoptions}
\opt{A}{E$_{cell}$ remains constant because anodic compartment contains buffer}
\opt{B}{E$_{cell}$ increases by 0.036 V}
\opt{C}{pH of anodic solution increases by 0.6}
\opt{D}{E$_{cell}$ decreases by 0.068 V}
\end{qoptions}

\qn{10} Which of the following statement(s) is/are correct ?
\begin{qoptions}
\opt{A}{In enzyme catalysed reactions, the cavities on enzyme ``Maltase'' and structure of molecules of Maltose have complimentary shapes for conversion of Maltase to glucose.}
\opt{B}{Most of the enzyme catalysed reactions are elementary.}
\opt{C}{Surface tension of lyophilic sols is less than the surface tension of dispersion medium.}
\opt{D}{Colloidal sols show lower magnitude of colligative properties, as compared to true solutions.}
\end{qoptions}

\sect{Section-II(i)}

\paraq{1}{2}
Treatment of compound \textbf{O} with KMnO$_4$/H$^+$ gave \textbf{P}, which on heating with ammonia gave \textbf{Q}. The compound \textbf{Q} on treatment with Br$_2$/NaOH produced \textbf{R}. On strong heating, \textbf{Q} gave \textbf{S}, which on further treatmenet with ethyl 2-bromopropanoate in the presence of KOH following by acidification, gave a compound T.
\qfig{p1_chem_s3_para1_2_1}

\qn{1} Find out total number of alkyl halide. Which react with \textbf{S} to give 1$^\circ$ \textbf{amine} as a products
\qfig{p1_chem_s3_q01_1}

\qn{2} What is degree of unsaturation of compound \textbf{P}

\paraq{3}{4}
In metallargy, self reduction is used for sulphides of Pb, Cu and Hg. The gas (A) released during this process when reacts with Zn + dil. H$_2$SO$_4$ produces another gas (B) which has characteristic rotten egg smell. Gas (B) when reacts with N, N-dimethyl-p-phenylene diamine in presence of acidic solution of Fe$^{+3}$(aq) gives a blue color dye (C).

\qn{3} Number of `N' atom present in `C' is/are ?

\qn{4} When gas `B' is passed through aqueous solution of `A' a yellowish white ppt is obtained. What is oxidation state of `S' in ppt.

\paraq{5}{6}
The mixing of different gases by frequent colllision is called diffusion and the process of escaping gas molecules through tiny orifice in vaccum is called effusion.

\qn{5} The reaction P$_{(g)}\to2$Q$_{(g)}+$R$_{(g)}$ follows I$^{\text{st}}$ order kinetics having rate constant 0.023 sec$^{-1}$. 1 mole compound P is taken in an evacuted sealed flask. After `t' time in the flask a hole is generated and the gases effuses according to Graham's law. The initially effusing gaseous mixture contains double the moles of Q than P. Determine value of t (in sec) if gram molecular masses of P, Q and R are 24, 6, 9 respectively.\\
\relax[Given : log3/2 = 0.176, ln 10 = 2.3]

\qn{6} The mass \% of R in the mixture coming out initially ? (Give nearest integer as your answer).

\qfullbegin
\sect{Section-II(ii)}

\qn{7} How many of following reactions will give racemic products :
\qfig{p1_chem_s4_q07_1}
\qfullsplit
\qfig{p1_chem_s4_q07_2}
\qfigi{p1_chem_s4_q07_3}\ $\xrightarrow[\text{2. H}_2\text{O}]{\text{1. CH}_3\text{MgBr}}$
\qfig{p1_chem_s4_q07_4}
\qfig{p1_chem_s4_q07_5}
\qfig{p1_chem_s4_q07_6}
\qfullend

\qn{8} How many stable conformer of compound
\qfigi{p1_chem_s4_q08_1}
possible which has nonzero dipole moment :

\qn{9} Calculate $\Delta$G for the given process :\\
X($\ell$, 1 atm, 300 K) $\longrightarrow$ X(g, 1 atm, 300 K)\\
Given that $\Delta$H$_{vap}$ at 500 K = 50 kJ / mol, C$_P$(X, $\ell$) = 100 J K$^{-1}$ mol$^{-1}$, C$_P$(X, g) = 58 J K$^{-1}$ mol$^{-1}$\\
Normal boiling point of liquid X is 500 K.\\
\relax[Use : $\ln\dfrac53=0.5$]\\
\textit{(give your answer nearest integer)}

\partsec{3}{MATHEMATICS}
\sect{Section-I(i)}

\qn{1} For a complex number $z=x+iy$ ; $\left(i=\sqrt{-1}\right)$ it satisfy $|z-1|+|z+3|\le8$ then all possible values of $|z-4|$ belongs to .........
\begin{qoptions}
\opt{A}{[1, 9]}
\opt{B}{[0, 4]}
\opt{C}{[2, 10]}
\opt{D}{[4, 8]}
\end{qoptions}

\qn{2} Image of the point A(1, 2, $-1$) with respect to the line $\ell:\dfrac{x-2}{3}=\dfrac{4y-4}{18}=\dfrac{2z-6}{18}$ as a mirror is A$'$ then distance of A$'$ from the plane $2x+3y+6z=16$, measured parallel to $\ell$ is -
\begin{qoptions}
\opt{A}{1}
\opt{B}{2}
\opt{C}{3}
\opt{D}{4}
\end{qoptions}

\qn{3} A four digit number is called a doublet if any of its digit is the same as only one neighbour for example 1221 is a doublet but 1222 is not, then number of such doublets are
\begin{qoptions}
\opt{A}{2349}
\opt{B}{2268}
\opt{C}{2259}
\opt{D}{2277}
\end{qoptions}

\qn{4} If A and B are independent events such that the probability that they both occur simultaneously is $\dfrac18$ and the probability that neither of them will occur is $\dfrac38$, then which of the following is always true ?
\begin{qoptions}
\opt{A}{the probability that event A occurs is $\dfrac12$}
\opt{B}{the probability that event B occurs is $\dfrac14$}
\opt{C}{the probability that event A occurs is less than $\dfrac34$}
\opt{D}{the probability that event B occurs is more than $\dfrac13$}
\end{qoptions}

\sect{Section-I(ii)}

\qn{5} Let $n\times n$ matrix is A$_n$ whose all elements are n except diagonal elements and given that a$_{ii}$ = i, for i = 1, 2, 3, ... n; then
\begin{qoptions}
\opt{A}{$|A_n|=n|A_{n-1}|$}
\opt{B}{$|A_n|=-n|A_{n-1}|$}
\opt{C}{$|A_n|=(-1)^n\,n!$}
\opt{D}{$|A_n|=(-1)^{n-1}\,n!$}
\end{qoptions}

\qn{6} Let $f(x)=\dfrac{-x+\sin x-[-x]+|\{-x\}-\sin x|}{2}$ then (where $\{\cdot\}$ denotes fractional part of x ; $[\cdot]$ denotes greatest integer $\le$ x)
\begin{qoptions}
\opt{A}{$f(x)$ is discontinuous at 7 points in $x\in(0,7)$}
\opt{B}{$f(x)$ is non-differentiable at 10 points in $x\in(0,7)$}
\opt{C}{$f(x)\ge0$ for $\forall\ x\in R$}
\opt{D}{$\displaystyle\lim_{x\to5^+}f(x)=1$}
\end{qoptions}

\qn{7} Which of the following statement is (are) true ?
\begin{qoptions}
\opt{A}{If A is a diagonal matrix of order $3\times3$, is commutative with every square matrix of order $3\times3$ under multiplication and trace(A) = 12 then $|A|=64$}
\opt{B}{Let each elements a$_{ij}$ of matrix A of order $3\times3$ ($|A|\ne0$) is multiplied by $\lambda^{i-j}$ and $|B|$ is resulting determinant then $|A|=|B|$}
\opt{C}{If A is square matrix of order 3 such that $|A|=2$ then $\left|\left(\mathrm{adj}(A^{-1})\right)^{-1}\right|=\dfrac14$}
\opt{D}{If A is a square matrix such that $A^2=A$ then $(I+A)^8=I+255A$}
\end{qoptions}

\qn{8} Let $f(x)$ is a non-zero differentiable function satisfying the relation $f(x)+f(x^2)+f(x^3)=f(x+x^2+x^3)\ \forall\ x\in R$, then which of the following is/are true ?
\begin{qoptions}
\opt{A}{$f\,'(-1)=f\,'(3)$}
\opt{B}{$f\,'(-1)=f\,'(1)$}
\opt{C}{$f\,'(1)=f\,'(3)$}
\opt{D}{$3f(-1)=f(3)$}
\end{qoptions}

\qn{9} The curve $y=\ell x^2+\ell x+\dfrac1{48}$ and $x=\ell y^2+\ell y+\dfrac1{48}$ ($\ell\ne0$) touch each other for
\begin{qoptions}
\opt{A}{two rational values of $\ell$}
\opt{B}{two irrational values of $\ell$}
\opt{C}{four rational values of $\ell$}
\opt{D}{four irrational values of $\ell$}
\end{qoptions}

\qn{10} Solution of the differential equation\\
$x(y\ln x-y\ln y-x-y)\,dy=y(x\ln y-x\ln x-y-x)\,dx$;\\
$y(1)=e$ is $ay\ln\left(\dfrac xy\right)=bx\ln\left(\dfrac yx\right)-ce+d$ where a, b, c and d are integers. Which of the following are correct.
\begin{qoptions}
\opt{A}{a + b = 2}
\opt{B}{ab + c + d = 1}
\opt{C}{ab + bc + cd + ad + bd + ac = 0}
\opt{D}{a + b + c + d = 2}
\end{qoptions}

\sect{Section-II(i)}

\paraq{1}{2}
Two sides of $\Delta$ABC, AB and AC are given by $2x+3y=7$ and $7x-3y=11$ respectively. Also internal angle bisector of $\angle$ABC is given by $2x+y+3=0$ then

\qn{1} Coordinate of C in $\Delta$ABC is C(K$_1$, K$_2$) then $\dfrac{94(K_1-K_2)}{15}$ is ..............

\qn{2} Slope of tangent drawn to $x^2=4y$ from point of intersection of line AB and AC is

\paraq{3}{4}
Given hyperbola H : $x^2-y^2-6x+2y+5=0$ having center `C' and tangent is drawn at point P $\left(3+\sqrt6,\ \sqrt3+1\right)$ to the H is intersect the Asymptotes of H at A and B then

\qn{3} Area of $\Delta$ABC is .....

\qn{4} If locus of center of the circle circumscribing $\Delta$ABC is $px^2+qy^2+rx+2y+t=0$ (where p, q, r, t are constant) then p + q $-$ r + t = ...

\paraq{5}{6}
A plane $\pi$ passes through the point (1, 1, 1) and is parallel to the vectors $\vec a=-\hat\imath+\hat\jmath$ and $\vec b=\hat\imath-\hat k$ and $\pi$ meets the co-ordinate axes in A, B and C forming the tetrahedron OABC where O is origin.

\qn{5} Sum of area of all lateral surfaces of tetrahedron is K$\left(3+\sqrt3\right)$ square unit then K is \underline{\hspace{1cm}}

\qn{6} Sum of square of distances from vertex to opposite faces of tetrahedron is \underline{\hspace{1cm}} unit.

\sect{Section-II(ii)}

\qn{7} Let $f(x)=x^3-2x^2-4x+1$ then number of different real solution of $f(f(x))=0$.

\qn{8} If $\displaystyle\sum_{n=1}^{\infty}\tan^{-1}\left(\dfrac{1}{n^2-9n+\frac{39}{2}}\right)=2\pi-\tan^{-1}\left(\dfrac ab\right)$ (where a and b are relatively prime number) then value of $|a+b|$ is .......

\qn{9} Let a, b, c are distinct positive real numbers such that $a+b+c\ge4$ then minimum value of $\dfrac{a^3}{(a-b)(a-c)}+\dfrac{b^3}{(b-c)(b-a)}+\dfrac{c^3}{(c-a)(c-b)}$ is ....

\end{multicols}
\end{document}
